\section{Schur's lemma, and Maschke's theorem}
Let \(G\) be a finite group and let \(k\) be a field. A \(k\)-representation
of \(G\) is a pair \((M,\rho)\), where \(M\) is a \(k\)-vector space and
\[
    \rho:G\longrightarrow \operatorname{GL}(M)
\]
is a group homomorphism.

\begin{proposition}
    Let \((M,\rho)\) be a \(k\)-representation of \(G\). Then \(\rho\) extends
    uniquely by \(k\)-linearity to a unital \(k\)-algebra homomorphism
    \[
        \widetilde{\rho}:kG\longrightarrow \operatorname{End}_k(M).
    \]
    Explicitly, if
    \[
        x=\sum_{g\in G}a_g g\in kG,
    \]
    then
    \[
        \widetilde{\rho}(x)
        =
        \sum_{g\in G}a_g \rho(g).
    \]
\end{proposition}

\begin{proof}
    Since \(kG\) is the free \(k\)-vector space with basis \(G\), the map
    \(\rho:G\to \operatorname{GL}(M)\subseteq \operatorname{End}_k(M)\) extends
    uniquely to a \(k\)-linear map
    \[
        \widetilde{\rho}:kG\to \operatorname{End}_k(M).
    \]
    We need to check that \(\widetilde{\rho}\) is multiplicative.

    Let
    \[
        x=\sum_{g\in G}a_g g,
        \qquad
        y=\sum_{h\in G}b_h h.
    \]
    Then
    \[
        xy
        =
        \sum_{g,h\in G}a_gb_h gh.
    \]
    Therefore
    \[
        \begin{aligned}
            \widetilde{\rho}(xy)
             & =
            \sum_{g,h\in G}a_gb_h \rho(gh)       \\
             & =
            \sum_{g,h\in G}a_gb_h \rho(g)\rho(h) \\
             & =
            \left(\sum_{g\in G}a_g\rho(g)\right)
            \left(\sum_{h\in G}b_h\rho(h)\right) \\
             & =
            \widetilde{\rho}(x)\widetilde{\rho}(y).
        \end{aligned}
    \]
    Moreover,
    \[
        \widetilde{\rho}(1_G)=\rho(1_G)=\operatorname{id}_M.
    \]
    Hence \(\widetilde{\rho}\) is a unital \(k\)-algebra homomorphism.
\end{proof}

Conversely, suppose we are given a unital \(k\)-algebra homomorphism
\[
    \rho:kG\longrightarrow \operatorname{End}_k(M).
\]
Then its restriction to \(G\subseteq kG\) gives a group homomorphism
\[
    G\longrightarrow \operatorname{GL}(M).
\]
Indeed, for \(g\in G\), the inverse of \(g\) in \(kG\) is \(g^{-1}\), so
\[
    \rho(g)\rho(g^{-1})
    =
    \rho(gg^{-1})
    =
    \rho(1)
    =
    \operatorname{id}_M.
\]
Thus \(\rho(g)\in \operatorname{GL}(M)\).

\begin{proposition}
    Giving a \(k\)-representation of \(G\) on \(M\) is equivalent to giving a
    left \(kG\)-module structure on \(M\).
\end{proposition}

\begin{proof}
    Let
    \[
        \rho:kG\longrightarrow \operatorname{End}_k(M)
    \]
    be the algebra homomorphism associated to a representation. Define
    \[
        kG\times M\longrightarrow M,
        \qquad
        (x,m)\longmapsto x\cdot m:=\rho(x)(m).
    \]
    Since \(\rho\) is a unital algebra homomorphism, we have
    \[
        (xy)\cdot m
        =
        x\cdot (y\cdot m),
        \qquad
        1\cdot m=m.
    \]
    Thus \(M\) becomes a left \(kG\)-module.

    Conversely, if \(M\) is a left \(kG\)-module, then every \(g\in G\) acts on
    \(M\) by an invertible \(k\)-linear map, with inverse given by the action of
    \(g^{-1}\). Hence we get a group homomorphism
    \[
        G\longrightarrow \operatorname{GL}(M),
        \qquad
        g\longmapsto (m\mapsto g\cdot m).
    \]
    These two constructions are inverse to each other.
\end{proof}

Therefore one has an equivalence of categories
\[
    \operatorname{Rep}_k(G)
    \simeq
    kG\text{-}\operatorname{Mod}.
\]
Under this equivalence, morphisms of representations are precisely
\(kG\)-module homomorphisms.

\begin{definition}[Regular representation]
    The \emph{left regular representation} of \(G\) is the representation on
    \(kG\) given by left multiplication:
    \[
        \rho:G\longrightarrow \operatorname{GL}(kG),
        \qquad
        \rho(g)(x)=gx.
    \]
    Equivalently, it is the \(kG\)-module \(kG\) acting on itself by left
    multiplication.
\end{definition}

The associated algebra homomorphism is
\[
    \rho:kG\longrightarrow \operatorname{End}_k(kG),
    \qquad
    \rho(a)(x)=ax.
\]

\begin{lemma}
    The algebra homomorphism
    \[
        \rho:kG\longrightarrow \operatorname{End}_k(kG)
    \]
    associated with the regular representation is injective. In particular, the
    regular representation is faithful.
\end{lemma}

\begin{proof}
    Suppose \(a\in kG\) satisfies
    \[
        \rho(a)=0.
    \]
    Then \(\rho(a)\) is the zero endomorphism of \(kG\). Applying it to the identity
    element \(1\in kG\), we get
    \[
        0=\rho(a)(1)=a\cdot 1=a.
    \]
    Therefore \(a=0\), and hence \(\rho\) is injective.
\end{proof}

\begin{lemma}[Schur's lemma]
    Let \((S,\rho)\) and \((S',\rho')\) be two irreducible \(k\) representations of
    \(G\). Equivalently, let \(S\) and \(S'\) be simple left \(kG\)-modules. Then:

    \begin{enumerate}
        \item If \(S\) and \(S'\) are not isomorphic, then
              \[
                  \operatorname{Hom}_{kG}(S,S')=0.
              \]

        \item The endomorphism algebra
              \[
                  \operatorname{End}_{kG}(S)
                  =
                  \operatorname{Hom}_{kG}(S,S)
              \]
              is a division \(k\)-algebra.

        \item If \(k=\overline{k}\), then
              \[
                  \operatorname{End}_{kG}(S)=k\operatorname{id}_S.
              \]
    \end{enumerate}
\end{lemma}

\begin{proof}
    Let
    \[
        f:S\longrightarrow S'
    \]
    be a \(kG\)-module homomorphism. Then \(\ker f\) is a submodule of \(S\), and
    \(\operatorname{Im} f\) is a submodule of \(S'\). Since \(S\) and \(S'\) are
    simple, we have
    \[
        \ker f=0 \quad \text{or} \quad \ker f=S,
    \]
    and
    \[
        \operatorname{Im} f=0 \quad \text{or} \quad \operatorname{Im} f=S'.
    \]
    If \(f\neq 0\), then \(\ker f\neq S\) and \(\operatorname{Im} f\neq 0\). Hence
    \[
        \ker f=0,
        \qquad
        \operatorname{Im} f=S'.
    \]
    Thus every nonzero \(kG\)-module homomorphism
    \[
        f:S\to S'
    \]
    is an isomorphism. Therefore, if \(S\not\cong S'\), no nonzero homomorphism can
    exist, and so
    \[
        \operatorname{Hom}_{kG}(S,S')=0.
    \]

    Taking \(S'=S\), we see that every nonzero element of
    \(\operatorname{End}_{kG}(S)\) is an automorphism of \(S\). Hence
    \[
        \operatorname{End}_{kG}(S)
    \]
    is a division \(k\)-algebra.

    Finally, assume that \(k=\overline{k}\). Let
    \[
        f\in \operatorname{End}_{kG}(S).
    \]
    Since \(S\) is finite-dimensional over the algebraically closed field \(k\),
    the \(k\)-linear map \(f:S\to S\) has an eigenvalue \(\lambda\in k\). Thus
    \[
        f-\lambda\operatorname{id}_S
    \]
    is not invertible as a \(k\)-linear map. But it is still a \(kG\)-module
    endomorphism of \(S\). Since \(\operatorname{End}_{kG}(S)\) is a division
    algebra, every nonzero element is invertible. Therefore
    \[
        f-\lambda\operatorname{id}_S=0.
    \]
    Hence
    \[
        f=\lambda\operatorname{id}_S.
    \]
    So
    \[
        \operatorname{End}_{kG}(S)=k\operatorname{id}_S.
    \]
\end{proof}

\begin{corollary}
    Let \(k=\overline{k}\), and let
    \[
        \rho:G\longrightarrow \operatorname{GL}_n(k)
    \]
    be an irreducible representation of \(G\). If \(M\in \operatorname{GL}_n(k)\)
    satisfies
    \[
        M^{-1}\rho(g)M=\rho(g)
    \]
    for all \(g\in G\), then \(M\) is a scalar matrix.
\end{corollary}

\begin{proof}
    The condition
    \[
        M^{-1}\rho(g)M=\rho(g)
    \]
    is equivalent to
    \[
        \rho(g)M=M\rho(g)
    \]
    for all \(g\in G\). Thus \(M\) is an endomorphism of the representation
    \((k^n,\rho)\), that is,
    \[
        M\in \operatorname{End}_{kG}(k^n).
    \]
    Since the representation is irreducible and \(k\) is algebraically closed,
    Schur's lemma gives
    \[
        \operatorname{End}_{kG}(k^n)=k\operatorname{id}_{k^n}.
    \]
    Therefore
    \[
        M=\lambda I_n
    \]
    for some \(\lambda\in k\).
\end{proof}

\begin{theorem}
    If
    \[
        \operatorname{char} k \mid |G|,
    \]
    then \(kG\) is not semisimple.
\end{theorem}

\begin{proof}
    Let
    \[
        \varepsilon:kG\longrightarrow k
    \]
    be the augmentation map, defined by
    \[
        \varepsilon\left(\sum_{g\in G}a_g g\right)
        =
        \sum_{g\in G}a_g.
    \]
    This is a homomorphism of left \(kG\)-modules, where \(k\) is the trivial
    \(kG\)-module.

    Assume, for contradiction, that \(kG\) is semisimple. Then the surjection
    \[
        \varepsilon:kG\longrightarrow k
    \]
    splits as a \(kG\)-module homomorphism. Hence there exists a \(kG\)-module
    homomorphism
    \[
        s:k\longrightarrow kG
    \]
    such that
    \[
        \varepsilon\circ s=\operatorname{id}_k.
    \]
    Let
    \[
        x=s(1)\in kG.
    \]
    Since \(s\) is a \(kG\)-module homomorphism and \(k\) is the trivial module, for
    every \(h\in G\) we have
    \[
        hx=x.
    \]
    Thus \(x\) is fixed by left multiplication by every element of \(G\).

    Write
    \[
        x=\sum_{g\in G}a_g g.
    \]
    The condition \(hx=x\) for all \(h\in G\) implies that all coefficients
    \(a_g\) are equal. Therefore
    \[
        x=a\sum_{g\in G}g
    \]
    for some \(a\in k\). Hence
    \[
        \varepsilon(x)
        =
        a|G|.
    \]
    But \(\operatorname{char} k\mid |G|\), so \(|G|=0\) in \(k\). Thus
    \[
        \varepsilon(x)=0.
    \]
    On the other hand,
    \[
        \varepsilon(x)
        =
        \varepsilon(s(1))
        =
        1,
    \]
    a contradiction. Therefore \(kG\) is not semisimple.
\end{proof}

\begin{theorem}[Maschke's theorem]
    If
    \[
        \operatorname{char} k\nmid |G|,
    \]
    then \(kG\) is semisimple.
\end{theorem}

\begin{proof}
    Let \(V\) be a left \(kG\)-module, and let
    \[
        W\subseteq V
    \]
    be a \(kG\)-submodule. We prove that \(W\) has a \(kG\)-submodule complement in
    \(V\).

    Let
    \[
        \iota:W\hookrightarrow V
    \]
    be the inclusion map. Since \(W\) is a \(kG\)-submodule, \(\iota\) is a
    \(kG\)-module homomorphism.

    As \(k\)-vector spaces, choose a complement \(W'\) of \(W\) in \(V\), so that
    \[
        V=W\oplus W'.
    \]
    Let
    \[
        p:V\longrightarrow W
    \]
    be the \(k\)-linear projection onto \(W\) with respect to this decomposition.
    Thus
    \[
        p\circ \iota=\operatorname{id}_W.
    \]
    The map \(p\) need not be a \(kG\)-module homomorphism. We now average it over
    the group.

    Let
    \[
        \rho:G\longrightarrow \operatorname{GL}(V)
    \]
    be the representation corresponding to the \(kG\)-module \(V\). Define
    \[
        \pi:V\longrightarrow W
    \]
    by
    \[
        \pi(v)
        =
        \frac{1}{|G|}
        \sum_{g\in G}
        \rho(g)^{-1}p(\rho(g)v).
    \]
    This is well-defined because \(|G|\neq 0\) in \(k\).

    We claim that \(\pi\) is a \(kG\)-module homomorphism. Let \(h\in G\). Then
    \[
        \begin{aligned}
            \rho(h)^{-1}\pi\rho(h)
             & =
            \rho(h)^{-1}
            \left(
            \frac{1}{|G|}
            \sum_{g\in G}
            \rho(g)^{-1}p\rho(g)
            \right)
            \rho(h)                                 \\
             & =
            \frac{1}{|G|}
            \sum_{g\in G}
            \rho(h)^{-1}\rho(g)^{-1}p\rho(g)\rho(h) \\
             & =
            \frac{1}{|G|}
            \sum_{g\in G}
            \rho(gh)^{-1}p\rho(gh).
        \end{aligned}
    \]
    Since \(g\mapsto gh\) is a permutation of \(G\), we get
    \[
        \rho(h)^{-1}\pi\rho(h)=\pi.
    \]
    Equivalently,
    \[
        \pi\rho(h)=\rho(h)\pi
    \]
    for all \(h\in G\). Hence \(\pi\) is a \(kG\)-module homomorphism.

    Next, for \(w\in W\), since \(W\) is a \(kG\)-submodule, we have
    \[
        \rho(g)w\in W
    \]
    for all \(g\in G\). Therefore
    \[
        p(\rho(g)w)=\rho(g)w.
    \]
    Thus
    \[
        \begin{aligned}
            \pi(w)
             & =
            \frac{1}{|G|}
            \sum_{g\in G}
            \rho(g)^{-1}p(\rho(g)w) \\
             & =
            \frac{1}{|G|}
            \sum_{g\in G}
            \rho(g)^{-1}\rho(g)w    \\
             & =
            \frac{1}{|G|}
            \sum_{g\in G}w          \\
             & =
            w.
        \end{aligned}
    \]
    Hence
    \[
        \pi\circ \iota=\operatorname{id}_W.
    \]

    Now set
    \[
        U=\ker \pi.
    \]
    Since \(\pi\) is a \(kG\)-module homomorphism, \(U\) is a \(kG\)-submodule of
    \(V\). We claim that
    \[
        V=W\oplus U.
    \]

    For every \(v\in V\), we have
    \[
        v=\pi(v)+\bigl(v-\pi(v)\bigr).
    \]
    Here \(\pi(v)\in W\), and
    \[
        \pi\bigl(v-\pi(v)\bigr)
        =
        \pi(v)-\pi^2(v)
        =
        \pi(v)-\pi(v)
        =
        0,
    \]
    because \(\pi\) is the identity on \(W\). Hence
    \[
        v-\pi(v)\in U.
    \]
    Thus \(V=W+U\).

    Finally, if \(w\in W\cap U\), then
    \[
        \pi(w)=w
    \]
    because \(\pi\) is the identity on \(W\), while
    \[
        \pi(w)=0
    \]
    because \(w\in U=\ker\pi\). Hence \(w=0\). Therefore
    \[
        W\cap U=0.
    \]
    So
    \[
        V=W\oplus U
    \]
    as \(kG\)-modules.

    Thus every \(kG\)-submodule of every \(kG\)-module has a complement. Therefore
    \(kG\) is semisimple.
\end{proof}
